\subsection{Overview}

This section presents the complete portfolio modeling framework developed through systematic literature calibration. The framework consists of four integrated modules:

\begin{enumerate}
    \item \textbf{Portfolio Composition:} Project categorization, profit margins, and optimal mix
    \item \textbf{Project Characteristics:} Portfolio size and BAC distributions
    \item \textbf{Cashflow Profiles:} S-curve models with Beta distributions
    \item \textbf{Revenue Plans:} Milestone-based payment structures
\end{enumerate}

Each module is documented with formal mathematical specifications, literature-calibrated parameters, and explicit assumptions. We adopt the \textit{Simplest Model Principle}: each choice represents minimal complexity for scientific validity, with extensions documented as future work.

\subsection{Portfolio Composition Model}

\subsubsection{Foundational Assumptions}

We assume all generated portfolios reflect \textit{optimal project selection} as characterized by \citet{khanzadi2018integrating}, who demonstrated that risk-adjusted return maximization (Sharpe ratio optimization) for Iranian EPC contractors yields:

\begin{equation}
\text{Optimal Composition: } 60\% \text{ domestic} + 40\% \text{ international}
\end{equation}

This composition balances:
\begin{itemize}
    \item \textbf{Domestic projects:} Lower profit margins (8--12\%), lower variance, regulatory stability
    \item \textbf{International projects:} Higher profit margins (12--16\%), higher variance, market-driven pricing
\end{itemize}

\textbf{Rationale:} Assuming optimal composition allows focus on operational dynamics (cashflow management, resource allocation) without solving the combinatorial project selection problem. This separation of concerns is standard in portfolio management literature \citep{archer1999integrated, cooper2001portfolio}.

\subsubsection{Profit Margin Distributions}

Profit margins are modeled as truncated normal distributions with category-specific parameters:

\textbf{Domestic Projects:}
\begin{equation}
\pi_i \sim \text{TN}(\mu = 0.10, \sigma = 0.012, a = 0.08, b = 0.12)
\end{equation}

Parameters reflect Iranian regulatory constraints: Budget Law mandates 8--12\% margins for government contracts, with 10\% as the standard rate \citep{khanzadi2018integrating}.

\textbf{International Projects:}
\begin{equation}
\pi_i \sim \text{TN}(\mu = 0.14, \sigma = 0.030, a = 0.12, b = 0.16)
\end{equation}

Higher mean and variance reflect market-driven competitive bidding without regulatory caps. Calibrated to \citet{khanzadi2018integrating} analysis of 127 Iranian EPC projects (2010--2016).

\textbf{Resulting distributions:}

\begin{table}[h]
\centering
\caption{Profit Margin Distributions by Category}
\label{tab:profit_margins}
\begin{tabular}{lcc}
\toprule
Percentile & Domestic & International \\
\midrule
P10 & 8.5\% & 10.6\% \\
P25 & 9.2\% & 12.1\% \\
Median & 10.0\% & 14.0\% \\
P75 & 10.8\% & 15.9\% \\
P90 & 11.5\% & 17.4\% \\
\bottomrule
\end{tabular}
\end{table}

\subsection{Portfolio Size and Project BAC Model}

\subsubsection{Portfolio Size Distribution}

Portfolio size $N$ (number of concurrent projects) is modeled as:

\begin{equation}
N \sim \text{DiscreteUniform}(N_{\min} = 8, N_{\max} = 18)
\end{equation}

\textbf{Calibration:} Based on empirical evidence from multiple sources:
\begin{itemize}
    \item \citet{merrow2011industrial}: Tier 1 EPC contractors manage 15--25 major projects; Tier 2 manage 8--15 projects
    \item CII (2019): Optimal diversification zone is 10--20 projects
    \item \citet{khanzadi2018integrating}: Iranian Tier 1 contractors manage 8--15 projects
\end{itemize}

Range captures Tier 1--2 Iranian EPC contractors in oil \& gas sector.

\subsubsection{Project BAC Distribution}

Project Budget at Completion (BAC) follows a lognormal distribution:

\begin{equation}
\log(\text{BAC}_i) \sim \mathcal{N}(\mu_{\log} = 5.01, \sigma_{\log} = 0.69)
\end{equation}

This yields:
\begin{itemize}
    \item Median BAC: \$150M
    \item Mean BAC: \$280M (right-skewed)
    \item Range: \$50M -- \$5B+
\end{itemize}

\textbf{Calibration:} \citet{merrow2011industrial} IPA database of 318 industrial megaprojects. The lognormal distribution captures the empirical pattern: most projects are mid-sized (\$100M--\$500M), with a long tail of megaprojects (\$1B+).

\textbf{Size categories:}
\begin{itemize}
    \item Small: \$50M -- \$200M (60\% of projects)
    \item Mid-size: \$200M -- \$1B (30\% of projects)
    \item Mega: $>$\$1B (10\% of projects)
\end{itemize}

\subsection{S-Curve Cashflow Model}

\subsubsection{Mathematical Formulation}

Planned cumulative spending for project $i$ at normalized time $\tau \in [0,1]$ follows a Beta distribution:

\begin{equation}
S_i(\tau) = \text{BAC}_i \times I_\tau(\alpha_i, \beta_i)
\end{equation}

where $I_\tau(\alpha, \beta)$ is the regularized incomplete beta function:

\begin{equation}
I_\tau(\alpha, \beta) = \frac{\int_0^\tau t^{\alpha-1}(1-t)^{\beta-1} dt}{B(\alpha, \beta)}
\end{equation}

Incremental spending rate (cashflow velocity):

\begin{equation}
c_i(\tau) = \frac{dS_i}{d\tau} = \text{BAC}_i \times \frac{\tau^{\alpha_i-1}(1-\tau)^{\beta_i-1}}{B(\alpha_i, \beta_i)}
\end{equation}

\textbf{Peak spending time:}

\begin{equation}
\tau^* = \frac{\alpha_i - 1}{\alpha_i + \beta_i - 2}
\end{equation}

\subsubsection{Parameter Calibration}

S-curve shape parameters $(\alpha_i, \beta_i)$ vary systematically with project characteristics. Base parameterization from \citet{cioffi2005tool}:

\begin{align}
\alpha_{\text{base}} &= 2.5 \\
\beta_{\text{base}} &= 2.0
\end{align}

This yields $\tau^* = 0.6$ (peak spending at 60\% of duration), matching empirical observations for typical EPC projects \citep{kenley1986construction}.

\textbf{Project-specific adjustments:}

\begin{equation}
\alpha_i = \alpha_{\text{base}} \times \left(\frac{\text{BAC}_i}{\text{BAC}_{\text{median}}}\right)^{0.15} \times m_{\text{category}} \times m_{\text{risk}} \times \epsilon_\alpha
\end{equation}

\begin{equation}
\beta_i = \beta_{\text{base}} \times \left(\frac{D_i}{D_{\text{median}}}\right)^{0.12} \times m_{\text{category}} \times m_{\text{risk}} \times \epsilon_\beta
\end{equation}

where:
\begin{itemize}
    \item Size effect: Larger projects have higher $\alpha$ (slower mobilization)
    \item Duration effect: Longer projects have higher $\beta$ (extended close-out)
    \item Category multipliers: $m_{\text{domestic}} = 0.9$, $m_{\text{international}} = 1.1$
    \item Risk multipliers: $m_{\text{low}} = 1.0$, $m_{\text{high}} = 0.8$ (for $\alpha$), $1.2$ (for $\beta$)
    \item Stochastic residual: $\epsilon_\alpha, \epsilon_\beta \sim \text{LogNormal}(0, 0.10)$
\end{itemize}

Parameters clipped to empirically validated bounds: $\alpha_i \in [1.2, 5.0]$, $\beta_i \in [1.2, 5.5]$ \citep{kenley1986construction}.

\textbf{Resulting parameter distributions by project type:}

\begin{table}[h]
\centering
\caption{S-Curve Parameters by Project Category}
\label{tab:scurve_params}
\begin{tabular}{llccc}
\toprule
Category & Risk & $\alpha$ & $\beta$ & Peak $\tau^*$ \\
\midrule
Domestic & Low & 2.00 & 1.80 & 0.556 \\
Domestic & High & 1.60 & 2.16 & 0.341 \\
International & Low & 2.50 & 2.16 & 0.600 \\
International & High & 2.00 & 2.59 & 0.435 \\
\bottomrule
\end{tabular}
\end{table}

\textbf{Interpretation:}
\begin{itemize}
    \item Domestic low-risk: Symmetric, slight front-load
    \item Domestic high-risk: Early peak, long tail (rework/disputes)
    \item International low-risk: Later peak, extended close-out
    \item International high-risk: Mid-peak, heavy tail (commissioning uncertainty)
\end{itemize}

\subsubsection{Project Duration Model}

Project duration $D_i$ (in months) correlates with BAC:

\begin{equation}
D_i = D_{\text{base}} \times \left(\frac{\text{BAC}_i}{\text{BAC}_{\text{ref}}}\right)^{0.35} \times \epsilon_D
\end{equation}

where:
\begin{itemize}
    \item $D_{\text{base}} = 24$ months (baseline for \$150M project)
    \item $\text{BAC}_{\text{ref}} = \$150M$ (reference project size)
    \item Exponent 0.35 reflects economies of scale (larger projects don't scale linearly in time)
    \item $\epsilon_D \sim \text{LogNormal}(0, 0.15)$ captures project-specific variation
\end{itemize}

Duration clipped to $[12, 60]$ months. Calibrated to \citet{merrow2011industrial} analysis showing median duration of 24 months for \$150M projects, with range 12--60 months across portfolio.

\subsection{Revenue Payment Plans Model}

\subsubsection{Milestone-Based Payment Structure}

We model milestone-based payment schedules, reflecting the dominant practice in EPC contracts (87\% of contracts per \citet{cui2010systems}). Total contract revenue:

\begin{equation}
R_i = \text{BAC}_i \times (1 + \pi_i)
\end{equation}

Revenue is released upon achievement of contractual milestones. Standard milestone structure:

\begin{table}[h]
\centering
\caption{Standard Milestone Payment Structure}
\label{tab:milestones}
\begin{tabular}{lcc}
\toprule
Milestone & Progress Threshold & Payment Weight \\
\midrule
Advance & 0\% (contract signing) & 10\% \\
Milestone 1 & 25\% & 20\% \\
Milestone 2 & 50\% & 30\% \\
Milestone 3 & 75\% & 30\% \\
Final & 100\% (completion) & 10\% \\
\bottomrule
\end{tabular}
\end{table}

\subsubsection{Payment Timing Linked to Actual Progress}

Critical modeling choice: milestone payments are triggered by \textit{actual achievement} of progress thresholds, not calendar dates. This links payment timing to contractor performance (SPI).

For milestone $m$ at planned progress $p_m$, payment is triggered when:

\begin{equation}
\text{progress}_t^i \geq p_m \quad \text{and} \quad \text{progress}_{t-1}^i < p_m
\end{equation}

If contractor is behind schedule ($\text{SPI} < 1.0$), milestone achievement is delayed, and so is payment. For a project with planned milestone at time $t_m^{\text{plan}}$ and actual SPI:

\begin{equation}
t_m^{\text{actual}} = \frac{t_m^{\text{plan}}}{\text{SPI}}
\end{equation}

\textbf{Example:}
\begin{itemize}
    \item Planned milestone: Month 12 (50\% progress)
    \item Actual SPI: 0.8 (20\% behind schedule)
    \item Actual milestone achievement: Month 12 / 0.8 = Month 15
    \item Payment received: Month 15 (not Month 12)
\end{itemize}

This creates direct coupling between execution performance and cash inflow timing, a key source of portfolio-level risk.

\subsubsection{Exclusions and Future Work}

The following elements are explicitly excluded:

\textbf{Retention Money:} We exclude Defects Liability Period (DLP) retention (typically 5--10\% withheld for 12--24 months post-completion). Rationale: DLP extends cashflow collection beyond strategic planning horizon; our focus is on active execution phase.

\textbf{Payment Delays:} We model milestone-triggered payments as instantaneous. Future work can add stochastic payment delays (e.g., 1--3 month processing time) to capture client payment risk.

\textbf{Alternative Payment Structures:} We exclude time-based payments, cost-plus contracts, and unit-price contracts. These represent natural extensions for broader applicability.

\subsection{Uncertainty Model: Stochastic Contractor Performance}

\subsubsection{Schedule Performance Index (SPI) Distribution}

SPI at completion follows a Beta distribution rescaled to $[0.3, 1.3]$:

\begin{equation}
\text{SPI}_{\text{final}}^i \sim 0.3 + 1.0 \times \text{Beta}(\alpha_{\text{SPI}} = 2.8, \beta_{\text{SPI}} = 2.2)
\end{equation}

\textbf{Calibration:} \citet{merrow2011industrial} analysis of 318 megaprojects:
\begin{itemize}
    \item Mean SPI: 0.81 (19\% behind schedule)
    \item Standard deviation: 0.19
    \item Left-skewed distribution (more underperformers)
\end{itemize}

Our parameterization yields mean SPI = 0.86 (slightly higher to account for action plan interventions in managed portfolios).

\textbf{Key percentiles:}

\begin{table}[h]
\centering
\caption{SPI Distribution Percentiles}
\label{tab:spi_dist}
\begin{tabular}{lcc}
\toprule
Percentile & SPI Value & Schedule Overrun \\
\midrule
P10 & 0.58 & 72\% \\
P25 & 0.72 & 39\% \\
P50 & 0.85 & 18\% \\
P75 & 0.98 & 2\% \\
P90 & 1.12 & -11\% (ahead) \\
\bottomrule
\end{tabular}
\end{table}

\subsubsection{Cost Performance Index (CPI) Distribution}

CPI at completion follows a Beta distribution rescaled to $[0.5, 1.2]$:

\begin{equation}
\text{CPI}_{\text{final}}^i \sim 0.5 + 0.7 \times \text{Beta}(\alpha_{\text{CPI}} = 3.2, \beta_{\text{CPI}} = 1.8)
\end{equation}

\textbf{Calibration:} \citet{love2016relationship} analysis of 276 construction projects:
\begin{itemize}
    \item Mean CPI: 0.87 (15\% over budget)
    \item Standard deviation: 0.14
    \item Correlation with SPI: $\rho = 0.68$
\end{itemize}

\subsubsection{SPI-CPI Correlation}

SPI and CPI are positively correlated ($\rho \approx 0.68$) due to shared root causes: projects behind schedule typically incur cost overruns from extended overhead and acceleration efforts \citep{love2016relationship}.

We enforce this correlation using a Gaussian copula:
\begin{enumerate}
    \item Sample $(u_{\text{SPI}}, u_{\text{CPI}})$ from bivariate normal with correlation $\rho = 0.68$
    \item Transform to uniform: $(v_{\text{SPI}}, v_{\text{CPI}}) = (\Phi(u_{\text{SPI}}), \Phi(u_{\text{CPI}}))$
    \item Apply inverse CDFs: $\text{SPI} = F_{\text{SPI}}^{-1}(v_{\text{SPI}})$, $\text{CPI} = F_{\text{CPI}}^{-1}(v_{\text{CPI}})$
\end{enumerate}

This preserves marginal distributions while enforcing target correlation.

\subsection{Summary of Modeling Framework}

The complete portfolio modeling framework integrates:

\begin{enumerate}
    \item \textbf{Composition:} 60/40 domestic/international mix with category-specific profit margins
    \item \textbf{Size:} 8--18 projects per portfolio, BAC lognormal (\$50M--\$5B+)
    \item \textbf{Cashflows:} Beta-distributed S-curves with project-specific shape parameters
    \item \textbf{Revenues:} Milestone-based payments linked to actual progress (SPI-adjusted)
    \item \textbf{Uncertainty:} Correlated SPI/CPI distributions calibrated to 318 megaprojects
\end{enumerate}

All parameters are literature-calibrated with explicit citations. The framework generates realistic portfolio instances suitable for RL training and baseline comparison. Complete parameter tables and validation results are provided in Appendix~\ref{app:calibration}.
